% ============================================================
%  GENERAL INTRODUCTION
% ============================================================

\chapter*{General Introduction}
\addcontentsline{toc}{chapter}{General Introduction}

The rapid spread of the Internet and mobile technologies has profoundly changed the way
people manage their daily health. Health applications installed on smartphones and
wearable devices now allow millions of users to monitor their physical condition, track
their habits, and access medical information at any time. According to the World Health
Organization, the use of mobile wireless technologies for health purposes not only gives
users access to information and services, but also contributes to positive changes in
health behaviours that help prevent both acute and chronic diseases. There are currently
more than seven billion mobile phone subscriptions worldwide, and in some countries
access to a mobile device is easier than access to clean water or electricity.

Despite this progress, the reality of daily healthcare remains fragmented in many
countries, including Tunisia. Patient medical records are not systematically digitalised,
and there is little or no structured communication between different healthcare facilities.
Patients are often forced to carry heavy paper folders containing a disorganised mix of
prescriptions, test results, and medical reports issued by different doctors at different
locations. This fragmentation leads at best to unnecessary repetition of examinations, and
at worst to dangerous situations where harmful combinations of medications go undetected
because no single practitioner has access to the full picture.

This problem takes on a more serious dimension when the patient is a pregnant woman.
Every health decision made during pregnancy affects not only the mother but also the
child she is carrying. A medication that is harmless in a healthy adult may cause serious
harm to a developing baby at a specific stage of pregnancy. Pregnant women in Tunisia
therefore face a dual challenge: navigating a fragmented healthcare system while managing
a level of medical complexity that demands precise and personalised information about
every treatment they take.

It is from this observation that the Pharmacie Nahla Noureddine, noticing a real and
recurring need among its patients, approached the development team and proposed the
creation of a dedicated digital health companion. The result is \textbf{Sahhty}, a word
meaning ``my health'' in Tunisian Arabic, a mobile application designed to give pregnant
women and all patients a single, unified tool to manage their health throughout their care
journey. The platform allows patients to record and monitor their health readings, follow
their pregnancy week by week, find and book appointments with nearby doctors, store and
share their medical documents securely, check the safety of medications during pregnancy,
and receive immediate alerts when something requires their attention. Doctors benefit from
a dedicated interface to manage their schedule, review authorised patient files, and respond
to appointment requests in real time.

This project was carried out as a final-year internship in partial fulfilment of the
Bachelor's degree in Computer Science at the Higher Institute of Sciences and Technology
of Hammam Sousse (ESSTHS), and was completed over a period of four months by a team
of two students working in close collaboration.

This report documents the full development of the Sahhty platform across six chapters.
The first chapter presents the project context, the problem that motivated it, a study of
existing solutions, and the development approach adopted by the team. The second chapter
defines the needs of the different users and the features the system must provide. The
third chapter describes how the system is structured and how its main interactions unfold.
The fourth chapter covers the machine learning engine used for health risk evaluation and
the medication safety module. The fifth chapter presents the implementation environment
and the main screens of the application. The sixth chapter reports on the testing and
validation activities carried out to confirm that the platform works correctly and meets
its objectives.

Together, these six chapters tell the complete story of a project that began with a real
health need and resulted in a functional and useful application for patients and healthcare
professionals alike.
